\section{Canaries with persistent state and delayed outcomes}\label{sec:C32:03}
A canary compares a new release with a baseline on a declared traffic population. For stateless requests, random assignment can often occur per request. For agents with persistent memory, assignment should usually occur at a level that prevents state contamination, such as user, repository, or session lineage. Switching versions turn by turn can make the new policy read memories written by the old policy and vice versa, confounding both treatment and safety interpretation.

Let $Z_i$ denote treatment assignment for independent cluster $i$, and $Y_i(z)$ its outcome under release $z$. The estimand can be mean task success, violation probability, completion time, or a vector of operational metrics. A release decision should specify acceptable changes before reading the results. A quality improvement cannot automatically compensate for an authorization regression. Delayed outcomes require waiting or explicitly handling censoring; early completion rates can favor the release that abandons difficult tasks quickly.

Shared retrieval indexes create interference even when sessions are assigned independently. If canary runs write new procedures into a global skill bank that baseline runs can retrieve, the baseline has been treated indirectly. Possible designs include isolated memory namespaces, frozen read-only shared evidence with separate write stores, or cluster-level assignment that includes the shared state. Each has a different cost and external-validity tradeoff. The experimental report must state which information crossed treatment boundaries.

Sequential monitoring creates repeated looks at the same data. Declaring success as soon as any metric crosses a conventional threshold inflates false positives. Use a predeclared stopping rule, a valid sequential procedure, or a fixed observation window appropriate to the estimand. Operational emergency stops are different: a serious violation can justify containment without first establishing statistical superiority or inferiority. The corresponding incident should be retained in the comparison rather than erased by restarting the experiment.

A proposed canary design uses four arms: baseline policy with baseline memory, new policy with baseline memory, new policy with migrated memory, and new policy with an isolated empty derived bank but the same raw evidence access. The first pair estimates a policy intervention at fixed retained state; the second and third isolate migration under the new policy; the third and fourth assess dependence on inherited derived state. These are conditional contrasts within the specified configurations, not a complete decomposition of all interactions. If the new policy cannot consume the baseline representation, that arm requires a declared compatibility adapter or controlled replay and cannot isolate a pure policy effect. No outcome is asserted. Additional cross-pairings can diagnose incompatibility, while real-effect exposure requires the corresponding operating authority and isolation.

